\documentclass[11pt,a4paper]{article}

\usepackage[T1]{fontenc}
\usepackage[utf8]{inputenc}
\usepackage{lmodern}
\usepackage[french]{babel}
\usepackage[a4paper,margin=2.3cm]{geometry}
\usepackage{microtype}
\usepackage{xcolor}
\usepackage{enumitem}
\usepackage{titlesec}
\usepackage{hyperref}
\usepackage{fancyvrb}
\usepackage[most]{tcolorbox}

\definecolor{ink}{HTML}{1c2330}
\definecolor{accent}{HTML}{2f6e4c}
\definecolor{accent2}{HTML}{8a5a1c}
\definecolor{soft}{HTML}{eef3ef}
\definecolor{softline}{HTML}{cfe0d4}
\definecolor{open}{HTML}{1c4e7a}
\definecolor{opensoft}{HTML}{eaf1f8}
\definecolor{theme}{HTML}{5a2f6e}
\definecolor{themesoft}{HTML}{f2ecf6}
\definecolor{codebg}{HTML}{f4f6f4}
\definecolor{codeline}{HTML}{d6e2d8}

\hypersetup{colorlinks=true, linkcolor=accent, urlcolor=accent2, pdftitle={Plan --- la suite}}

\titleformat{\section}{\Large\bfseries\color{ink}}{\thesection.}{0.5em}{}
\titlespacing{\section}{0pt}{1.2em}{0.4em}
\setlist[itemize]{leftmargin=1.2em,itemsep=2pt,topsep=2pt}

\newtcolorbox{constat}{enhanced,breakable,colback=soft,colframe=accent,boxrule=0.8pt,left=10pt,right=10pt,top=7pt,bottom=7pt,arc=3pt}
\newtcolorbox{openbox}{enhanced,breakable,colback=opensoft,colframe=open,boxrule=0.5pt,left=9pt,right=9pt,top=6pt,bottom=6pt,arc=2pt}
\newtcolorbox{phase}[1]{enhanced,breakable,colback=white,colframe=softline,boxrule=0.7pt,left=9pt,right=9pt,top=6pt,bottom=6pt,arc=2pt,fonttitle=\bfseries,coltitle=ink,title={#1}}
\DefineVerbatimEnvironment{code}{Verbatim}{fontsize=\small,frame=single,framerule=0.5pt,rulecolor=\color{codeline},formatcom=\color{ink}}

\title{\vspace{-1.5em}\textbf{\color{ink}\Huge Plan --- la suite}\\[0.3em]
\large\color{accent}Feuille de route après l'Acte 1 jouable\\[0.2em]
\normalsize\color{ink}\itshape Des piliers de la DIRECTION aux incréments jouables}
\author{}\date{}

\begin{document}
\maketitle
\vspace{-3em}

\section{Où on en est}
\textbf{En place :} éditeurs Mission / Campagne / Geoscape, runtime hub geoscape (carte $\to$ mission de région $\to$ retour $\to$ déverrouillage $\to$ \texttt{endWhen}), moteur de combat tour-par-tour (RNG, pods, vigilance, crackers, potions, brouillard), effet «~île~». \textbf{L'Acte 1 se joue en avant} : Opening $\to$ mission d'intro $\to$ geoscape $\to$ fin.

\smallskip
\textbf{Ce qui manque --- le cœur RPG.} Aujourd'hui : un \emph{bac à sable tactique + une coquille de campagne}. Presque tous les piliers de la DIRECTION ne sont pas encore dans le jeu : progression par choix, usure, mort à enjeu, roster persistant.

\begin{constat}
\textbf{Le constat directeur.} Ce qui transforme le projet en \textbf{le jeu de la DIRECTION}, c'est une boucle qui n'existe pas encore :\\[2pt]
\centering\textbf{roster persistant $\to$ déploiement $\to$ progression par choix $\to$ usure qui force la rotation $\to$ mort à enjeu.}\\[2pt]
\flushleft Tout s'accroche au \textbf{roster persistant} --- la fondation manquante (chaque mission définit ses unités ; la campagne ne transporte que les PV par nom).
\end{constat}

\section{Les phases}

\begin{phase}{Phase A --- Roster persistant + déploiement \textnormal{(\emph{fondation})}}
Un \textbf{escadron de campagne} : persos nommés, persistants, stockés dans l'état de partie. En mission lancée depuis le geoscape, le joueur \textbf{déploie son roster} sur des cases de départ (au lieu des unités pré-placées). Éditeur : marquer des \textbf{zones de déploiement}.\\
\textbf{Livrable :} lancer une région $\to$ choisir/déployer ses persos $\to$ jouer avec eux. \textbf{Débloque} B, C, D.
\end{phase}

\begin{phase}{Phase B --- Progression par perks \textnormal{(\emph{pilier «~décision~»}, B1/B3)}}
Grades \textbf{discrets} ; à chaque grade un \textbf{choix A/B irréversible} --- un nouveau \emph{verbe}, pas un plus gros chiffre. XP/puissance plafonnée. Écran de montée au camp.\\
\textbf{Livrable :} un perso monte d'un cran, on choisit A ou B, deux persos identiques divergent durablement.
\end{phase}

\begin{phase}{Phase C --- Stress + Fatigue $\to$ rotation \textnormal{(\emph{le fil rouge}, C3)}}
Deux jauges par perso : montent en mission, \textbf{récupèrent lentement au camp à slots limités}. Un perso trop usé devient \textbf{temporairement indisponible}.\\
\textbf{Livrable :} enchaîner les missions \textbf{force à faire tourner le roster}, donc à investir dans plusieurs builds.
\end{phase}

\begin{phase}{Phase D --- Mort \& statut spécial \textnormal{(\emph{l'enjeu}, C1/B9)}}
\textbf{KO vs mort définitive} (asymétrique) ; \textbf{statut spécial temporaire} des persos clés (protégés tant qu'ils servent l'histoire, mortels ensuite). Un mort \textbf{quitte le roster}.\\
\textbf{Livrable :} perdre un perso compte vraiment, sans bloquer le récit.
\end{phase}

\begin{phase}{Phase E --- Pression \& échec-fait-avancer \textnormal{(\emph{geoscape vivant}, A1/C4)}}
\textbf{Horloge / menace locale par région} ; un \textbf{échec change la carte} (région perdue, menace qui s'étend) au lieu de bloquer ; le \textbf{camp} devient le hub de récupération.\\
\textbf{Livrable :} tarder ou échouer \textbf{transforme le territoire}.
\end{phase}

\begin{phase}{Phase F --- Thème \& feel \textnormal{(\emph{polish})}}
\textbf{Déformation progressive de la carte} au fil de l'acte (thème «~tout se déforme~») ; trancher le \textbf{dosage RNG} (friction 1) manette en main ; \textbf{identité/customisation} des persos.\\
\textbf{Livrable :} le thème se \emph{voit} et se \emph{joue} ; le hasard du combat est calibré.
\end{phase}

\section{Dépendances}
\begin{code}
A (roster) -+-> B (perks)
            +-> C (stress/fatigue) --+
            +-> D (mort) ------------+--> E (pression / echec-fait-avancer)
                                         F (theme & feel) -- en dernier, transverse
\end{code}
\begin{itemize}
\item \textbf{A est bloquant} pour B, C, D.
\item \textbf{E} s'appuie sur D + C + le geoscape (déjà là).
\item \textbf{F} en dernier (ou par petites touches en continu).
\end{itemize}

\section{Décision encore ouverte}
\begin{openbox}
\textbf{Dosage RNG} (friction 1) : socle = RNG tempéré + curseur de variance, à trancher \textbf{en jouant} --- idéalement pendant B/C, quand il y a assez de matière.
\end{openbox}

\section{Cadence}
Travail \textbf{autonome, phase par phase} : implémentation $\to$ vérification Chromium (headless) $\to$ commit/push $\to$ aperçu. Un seul déploiement Pages à la fois (laisser chaque build finir avant de pousser le suivant).

\begin{center}\itshape Prochaine étape : Phase A --- Roster persistant + déploiement.\end{center}

\end{document}
